% 第1章贯穿图：V={1,2,3,4}，E={{1,2},{1,3},{2,3},{2,4},{3,4}}的三种画法。
\begin{tikzpicture}[panelmark/.style={font=\figurefont\fontsize{8}{10}\selectfont,text=FigureInk,anchor=north,inner sep=0pt,outer sep=0pt},
>=Stealth,
  v/.style={circle,draw=FigureStroke,fill=FigureBlueFill,line width=.8pt,minimum size=6.5mm,inner sep=0pt},
  e/.style={draw=FigureStroke,line width=.8pt},
  note/.style={font=\FigureNoteFont\fontsize{7.5}{9}\selectfont,text=FigureMuted}]
\begin{scope}[local bounding box=panelrow]
\begin{scope}[local bounding box=panela]
  \node[v] (1) at (-1.3,0) {$1$}; \node[v] (2) at (0,.9) {$2$};
  \node[v] (3) at (0,-.9) {$3$}; \node[v] (4) at (1.3,0) {$4$};
  \draw[e] (1)--(2) (1)--(3) (2)--(3) (2)--(4) (3)--(4);

\end{scope}
\begin{scope}[xshift=3.9cm,local bounding box=panelb]
  \node[v] (1) at (-1.35,0) {$1$}; \node[v] (2) at (-.45,0) {$2$};
  \node[v] (3) at (.45,0) {$3$}; \node[v] (4) at (1.35,0) {$4$};
  \draw[e] (1)--(2) (2)--(3) (3)--(4);
  \draw[e] (1) to[bend left=55] (3);
  \draw[e] (2) to[bend right=55] (4);

\end{scope}
\begin{scope}[xshift=7.6cm,local bounding box=panelc]
  \node[v] (1) at (-.8,.8) {$1$}; \node[v] (4) at (.8,.8) {$4$};
  \node[v] (3) at (-.8,-.8) {$3$}; \node[v] (2) at (.8,-.8) {$2$};
  \draw[e] (1)--(2) (1)--(3) (2)--(3) (2)--(4) (3)--(4);

\end{scope}
\end{scope}
% Subfigure marks: 8 pt, centred below each panel; common top 3 mm below this row.
\coordinate (panelrowbase) at ($(panelrow.south)+(0,-3mm)$);
\node[panelmark] at (panela.center |- panelrowbase) {(a)};
\node[panelmark] at (panelb.center |- panelrowbase) {(b)};
\node[panelmark] at (panelc.center |- panelrowbase) {(c)};
\end{tikzpicture}
